% Options for packages loaded elsewhere
\PassOptionsToPackage{unicode}{hyperref}
\PassOptionsToPackage{hyphens}{url}
\documentclass[
  a4paper,
]{ltjsarticle}
\usepackage{xcolor}
\usepackage[margin=25mm]{geometry}
\usepackage{amsmath,amssymb}
\setcounter{secnumdepth}{-\maxdimen} % remove section numbering
\usepackage{iftex}
\ifPDFTeX
  \usepackage[T1]{fontenc}
  \usepackage[utf8]{inputenc}
  \usepackage{textcomp} % provide euro and other symbols
\else % if luatex or xetex
  \usepackage{unicode-math} % this also loads fontspec
  \defaultfontfeatures{Scale=MatchLowercase}
  \defaultfontfeatures[\rmfamily]{Ligatures=TeX,Scale=1}
\fi
\usepackage{lmodern}
\ifPDFTeX\else
  % xetex/luatex font selection
\fi
% Use upquote if available, for straight quotes in verbatim environments
\IfFileExists{upquote.sty}{\usepackage{upquote}}{}
\IfFileExists{microtype.sty}{% use microtype if available
  \usepackage[]{microtype}
  \UseMicrotypeSet[protrusion]{basicmath} % disable protrusion for tt fonts
}{}
\makeatletter
\@ifundefined{KOMAClassName}{% if non-KOMA class
  \IfFileExists{parskip.sty}{%
    \usepackage{parskip}
  }{% else
    \setlength{\parindent}{0pt}
    \setlength{\parskip}{6pt plus 2pt minus 1pt}}
}{% if KOMA class
  \KOMAoptions{parskip=half}}
\makeatother
\usepackage{longtable,booktabs,array}
\newcounter{none} % for unnumbered tables
\usepackage{calc} % for calculating minipage widths
% Correct order of tables after \paragraph or \subparagraph
\usepackage{etoolbox}
\makeatletter
\patchcmd\longtable{\par}{\if@noskipsec\mbox{}\fi\par}{}{}
\makeatother
% Allow footnotes in longtable head/foot
\IfFileExists{footnotehyper.sty}{\usepackage{footnotehyper}}{\usepackage{footnote}}
\makesavenoteenv{longtable}
\usepackage{graphicx}
\makeatletter
\newsavebox\pandoc@box
\newcommand*\pandocbounded[1]{% scales image to fit in text height/width
  \sbox\pandoc@box{#1}%
  \Gscale@div\@tempa{\textheight}{\dimexpr\ht\pandoc@box+\dp\pandoc@box\relax}%
  \Gscale@div\@tempb{\linewidth}{\wd\pandoc@box}%
  \ifdim\@tempb\p@<\@tempa\p@\let\@tempa\@tempb\fi% select the smaller of both
  \ifdim\@tempa\p@<\p@\scalebox{\@tempa}{\usebox\pandoc@box}%
  \else\usebox{\pandoc@box}%
  \fi%
}
% Set default figure placement to htbp
\def\fps@figure{htbp}
\makeatother
\setlength{\emergencystretch}{3em} % prevent overfull lines
\providecommand{\tightlist}{%
  \setlength{\itemsep}{0pt}\setlength{\parskip}{0pt}}
% Glyph map for lualatex (newunicodechar). Maps symbols that may lack font glyphs.
\usepackage{newunicodechar}
\usepackage{amssymb}
\newunicodechar{∎}{\ensuremath{\blacksquare}}
\newunicodechar{✓}{\ensuremath{\checkmark}}
\newunicodechar{✗}{\ensuremath{\times}}
\newunicodechar{⟹}{\ensuremath{\Longrightarrow}}
\newunicodechar{⟺}{\ensuremath{\Longleftrightarrow}}
\newunicodechar{↪}{\ensuremath{\hookrightarrow}}
\newunicodechar{⌊}{\ensuremath{\lfloor}}
\newunicodechar{⌋}{\ensuremath{\rfloor}}
\newunicodechar{≃}{\ensuremath{\simeq}}
\newunicodechar{≈}{\ensuremath{\approx}}
\newunicodechar{≤}{\ensuremath{\leq}}
\newunicodechar{≥}{\ensuremath{\geq}}
\newunicodechar{≠}{\ensuremath{\neq}}
\newunicodechar{→}{\ensuremath{\rightarrow}}
\newunicodechar{←}{\ensuremath{\leftarrow}}
\newunicodechar{↔}{\ensuremath{\leftrightarrow}}
\newunicodechar{⊥}{\ensuremath{\perp}}
\newunicodechar{√}{\ensuremath{\surd}}
\newunicodechar{∑}{\ensuremath{\sum}}
\newunicodechar{∏}{\ensuremath{\prod}}
\newunicodechar{∞}{\ensuremath{\infty}}
\newunicodechar{∈}{\ensuremath{\in}}
\newunicodechar{∀}{\ensuremath{\forall}}
\newunicodechar{∣}{\ensuremath{\mid}}
\newunicodechar{γ}{\ensuremath{\gamma}}
\newunicodechar{λ}{\ensuremath{\lambda}}
\newunicodechar{μ}{\ensuremath{\mu}}
\newunicodechar{ρ}{\ensuremath{\rho}}
\newunicodechar{σ}{\ensuremath{\sigma}}
\newunicodechar{φ}{\ensuremath{\varphi}}
\newunicodechar{ψ}{\ensuremath{\psi}}
\newunicodechar{θ}{\ensuremath{\theta}}
\newunicodechar{Δ}{\ensuremath{\Delta}}
\newunicodechar{Π}{\ensuremath{\Pi}}
\newunicodechar{Λ}{\ensuremath{\Lambda}}
\newunicodechar{τ}{\ensuremath{\tau}}
\newunicodechar{ε}{\ensuremath{\varepsilon}}
\newunicodechar{ω}{\ensuremath{\omega}}
\newunicodechar{π}{\ensuremath{\pi}}
\newunicodechar{ν}{\ensuremath{\nu}}
\newunicodechar{±}{\ensuremath{\pm}}
\newunicodechar{×}{\ensuremath{\times}}
\newunicodechar{·}{\ensuremath{\cdot}}
\usepackage{bookmark}
\IfFileExists{xurl.sty}{\usepackage{xurl}}{} % add URL line breaks if available
\urlstyle{same}
\hypersetup{
  pdftitle={インフレーション曲線は床の線形不安定モードの指数増幅である------床ヤコビアンの最大固有値と非正規・時変補正},
  pdfauthor={Noriaki Kihara (WF System Co., Ltd.), ORCID 0009-0004-6753-4020},
  hidelinks,
  pdfcreator={LaTeX via pandoc}}

\title{インフレーション曲線は床の線形不安定モードの指数増幅である------床ヤコビアンの最大固有値と非正規・時変補正}
\author{Noriaki Kihara (WF System Co., Ltd.), ORCID 0009-0004-6753-4020}
\date{2026-09-12}

\begin{document}
\maketitle

\textbf{シリーズ:} 自己無撞着な関係波閉鎖系の基礎（論文6／全10本）\\
\textbf{著者:} 木原 範昭（WF System Co., Ltd.）　\textbf{日付:}
2026-09-12\\
\textbf{Version DOI:} 10.5281/zenodo.22729096\\
\textbf{Concept DOI:} 10.5281/zenodo.22729095\\
\textbf{ORCID:} 0009-0004-6753-4020\\
\textbf{Zenodo:} https://zenodo.org/records/22729096\\

\begin{center}\rule{0.5\linewidth}{0.5pt}\end{center}

\subsection{要旨}\label{ux8981ux65e8}

先行研究［1］が報告したインフレーション的急拡大（床平面外成分
\(H_\perp/H\)
の指数増大）を、微小シードの\textbf{線形不安定性}として解析的に導く。床の線形化写像（ヤコビアン
\(J\)：\(z\to\exp(\Delta\tau K(\varphi))z\) を床で線形化した
\(2M\times2M\) 実行列）の最大固有値 \(|\mu_{\max}|\)
が、\(H_\perp/H(t)\approx\text{seed}^2\,|\mu_{\max}|^{2t}\)、ランプ傾き
\(2\log_{10}|\mu_{\max}|\)、onset
\(\approx\ln(1/H_{\perp0})/(2\ln|\mu_{\max}|)=\ln(1/\text{seed})/\ln|\mu_{\max}|\)（\(H_{\perp0}=\text{seed}^2\)）を与える。固定床
\(J\) で微小 \(\delta\) を発展させた \(H_\perp\) 傾きは
\(2\log_{10}|\mu_{\max}|\) に厳密一致（\(N=6{:}0.2935\) vs 解析
\(0.2937\)）。\(|\mu_{\max}|\) は \(N\)
増で減少（\(1.402\to1.269\to1.185\)、\(N=6,8,10\)）＝インフレが遅くなり
onset が伸びる（実測 onset は \(N=6,8,10\) で \(64,87,133\)
歩）。実測ランプ傾きは床値より一貫して \(18\)--\(31\%\) 急で、これは床
\(J\) が強い\textbf{非正規行列}（非正規度
\(0.53\)--\(0.76\)）で、実ランプが \(K(\varphi)\)
の漸進変化に伴う\textbf{時変・非可換な時間順序積} \(\prod J(\varphi_t)\)
の増幅（同時スペクトル半径 \(\ge\)
個別スペクトル半径）であることに由来する（\textbf{追試Aで直接確認}：時変積の傾きが実測に一致・床凍結を18--31\%上回る）。床固有値はその下界（初期・凍結
rate）である。分類：導ける（形・機構・\(N\)
依存・onset）＋非正規補正の直接確認（追試A、閉形式のみ未決）。

\begin{center}\rule{0.5\linewidth}{0.5pt}\end{center}

\subsection{1. 課題}\label{ux8ab2ux984c}

論文1・2で床が不安定な相対平衡と確定した。「波の力学が解けたなら、インフレーション曲線も微小シードから解析的に導けるはず」------本論文はこれを、床ヤコビアンの線形不安定モードの指数増幅として定量的に示す。

\subsection{2.
線形不安定の予測式}\label{ux7ddaux5f62ux4e0dux5b89ux5b9aux306eux4e88ux6e2cux5f0f}

床 \(z_0\) の周りで写像を線形化した実ヤコビアン
\(J\in\mathbb R^{2M\times2M}\)（状態を \([\Re z;\Im z]\)
に埋め込む）を数値微分で構成し固有値 \(\mu\) を求める。床が不安定なら
\(|\mu_{\max}|>1\)。微小シード \(\delta\)
は不安定固有方向で増幅し、床平面直交成分の割合は

\[\frac{H_\perp}{H}(t)\approx(\text{seed})^2\,|\mu_{\max}|^{2t},\qquad
\text{ランプ傾き}=\frac{d\log_{10}(H_\perp/H)}{dt}=2\log_{10}|\mu_{\max}|,\qquad
\text{onset}\approx\frac{\ln(1/H_{\perp0})}{2\ln|\mu_{\max}|}=\frac{\ln(1/\text{seed})}{\ln|\mu_{\max}|}.\]

（onset は \(H_\perp/H\to1\)
の到達歩数。\(H_{\perp0}=\text{seed}^2\)。旧稿の
\(\ln(1/\text{seed})/(2\ln|\mu|)\) は係数誤りで、実測 onset
の約半分になる自己矛盾があったため訂正した。）

\subsection{\texorpdfstring{3.
検証（\(N=6,8,10\)）}{3. 検証（N=6,8,10）}}\label{ux691cux8a3cn6810}

{\def\LTcaptype{none} % do not increment counter
\begin{longtable}[]{@{}
  >{\raggedright\arraybackslash}p{(\linewidth - 18\tabcolsep) * \real{0.1000}}
  >{\raggedright\arraybackslash}p{(\linewidth - 18\tabcolsep) * \real{0.1000}}
  >{\raggedright\arraybackslash}p{(\linewidth - 18\tabcolsep) * \real{0.1000}}
  >{\raggedright\arraybackslash}p{(\linewidth - 18\tabcolsep) * \real{0.1000}}
  >{\raggedright\arraybackslash}p{(\linewidth - 18\tabcolsep) * \real{0.1000}}
  >{\raggedright\arraybackslash}p{(\linewidth - 18\tabcolsep) * \real{0.1000}}
  >{\raggedright\arraybackslash}p{(\linewidth - 18\tabcolsep) * \real{0.1000}}
  >{\raggedright\arraybackslash}p{(\linewidth - 18\tabcolsep) * \real{0.1000}}
  >{\raggedright\arraybackslash}p{(\linewidth - 18\tabcolsep) * \real{0.1000}}
  >{\raggedright\arraybackslash}p{(\linewidth - 18\tabcolsep) * \real{0.1000}}@{}}
\toprule\noalign{}
\begin{minipage}[b]{\linewidth}\raggedright
\(N\)
\end{minipage} & \begin{minipage}[b]{\linewidth}\raggedright
\(|\mu_{\max}|\)
\end{minipage} & \begin{minipage}[b]{\linewidth}\raggedright
\(\lambda=\ln|\mu|\)/step
\end{minipage} & \begin{minipage}[b]{\linewidth}\raggedright
不安定モード数
\end{minipage} & \begin{minipage}[b]{\linewidth}\raggedright
解析傾き \(2\log_{10}|\mu|\)
\end{minipage} & \begin{minipage}[b]{\linewidth}\raggedright
固定床\(J\)の\(\delta\)発展傾き
\end{minipage} & \begin{minipage}[b]{\linewidth}\raggedright
実測ランプ傾き
\end{minipage} & \begin{minipage}[b]{\linewidth}\raggedright
実測/床解析
\end{minipage} & \begin{minipage}[b]{\linewidth}\raggedright
床\(J\)非正規度
\end{minipage} & \begin{minipage}[b]{\linewidth}\raggedright
onset 解析/実測
\end{minipage} \\
\midrule\noalign{}
\endhead
\bottomrule\noalign{}
\endlastfoot
6 & 1.40234 & 0.3381 & 5 & 0.2937 & 0.2935 & 0.3460 & 1.178 & 0.761 & 86
/ 64 \\
8 & 1.26908 & 0.2383 & 6 & 0.2070 & 0.1988 & 0.2702 & 1.306 & 0.641 &
125 / 87 \\
10 & 1.18513 & 0.1699 & 8 & 0.1475 & 0.1271 & 0.1814 & 1.230 & 0.530 &
180 / 133 \\
\end{longtable}
}

\begin{itemize}
\tightlist
\item
  \textbf{線形不安定式は厳密に正しい}：固定床 \(J\) で微小 \(\delta\)
  を発展させた \(H_\perp\) 傾きが解析 \(2\log_{10}|\mu_{\max}|\)
  に一致（\(N=6{:}0.2935\) vs \(0.2937\)）。
\item
  \textbf{\(N\) 依存を再現}：\(|\mu_{\max}|\) は \(N\)
  増で減少（\(1.402\to1.269\to1.185\)）＝インフレが遅くなり onset
  が伸びる（実測 onset は \(N=6,8,10\) で \(64,87,133\) 歩、解析
  \(86,125,180\) 歩）。onset の閉形式 \(N\)
  依存則は今後の検討課題とする。
\item
  \textbf{onset} \(\approx\ln(1/\text{seed})/\ln|\mu|\)
  が桁・スケールで一致（解析 \(86/125/180\) vs 実測 \(64/87/133\)）。
\end{itemize}

\begin{figure}
\centering
\pandocbounded{\includegraphics[keepaspectratio,alt={図1: インフレ=線形不安定。左=床ヤコビアン\textbar μ\_max\textbar とonsetのN依存、中=解析ランプ傾き2log10\textbar μ\textbar(床)vs実測(非正規ギャップ)、右=N=6のH⊥/H実測ランプに解析傾きを重畳。}]{figs/p6_inflation_linear_instability_ja_fig1.png}}
\caption{図1:
インフレ=線形不安定。左=床ヤコビアン\textbar μ\_max\textbar とonsetのN依存、中=解析ランプ傾き2log10\textbar μ\textbar(床)vs実測(非正規ギャップ)、右=N=6のH⊥/H実測ランプに解析傾きを重畳。}
\end{figure}

\textbf{図1} インフレーション＝床の線形不安定（\(N\)
依存・傾き・重畳）。

\subsection{4.
非正規補正（実測が床値を超える分）}\label{ux975eux6b63ux898fux88dcux6b63ux5b9fux6e2cux304cux5e8aux5024ux3092ux8d85ux3048ux308bux5206}

実測ランプ傾きは床値より一貫して \(18\)--\(31\%\) 急である（実測/床解析
\(=1.18\)--\(1.31\)）。

\begin{itemize}
\tightlist
\item
  床 \(J\) は強い\textbf{非正規行列}（非正規度
  \(\|JJ^{\mathsf T}-J^{\mathsf T}J\|/\|J^{\mathsf T}J\|=0.53\)--\(0.76\)）。
\item
  各点で \(J\)
  を凍結したときのスペクトル半径はランプ上でほぼ一定（\(N=6\) で
  \(1.402\)）だが、実ランプは\textbf{持続的に}これを上回る。凍結非正規性は一時的（transient）増幅を与えるに過ぎず、持続分は
  \(K(\varphi)\)
  が漸進変化する\textbf{時変・非可換なヤコビアンの時間順序積}
  \(\prod_t J(\varphi_t)\) の増幅（同時スペクトル半径 \(\ge\)
  個別スペクトル半径）が担う。\textbf{この機構は直接検証で確定した（追試A）}：実軌道上の時変ヤコビアン積
  \(\prod_t J(z_t)\) で \(\delta\) を発展させた \(H_\perp\)
  傾きが実測ランプに一致し（時変/実測 =
  0.999／1.004／1.019、\(N=6,8,10\)）、床凍結 \(J_0\)
  を18--31\%上回る（時変/床 =
  1.177／1.311／1.253＝実測の超過と定量一致）。床固有値はその下界（初期・凍結
  rate）である。したがって §3 で解析 onset（床 \(|\mu_{\max}|\)
  基準）が実測を一貫して上回る（\(86/125/180\) vs 実測
  \(64/87/133\)）のも同じ非正規機構の帰結である------床固有値は成長率の下界ゆえ、床基準の
  onset は上界を与え、実ランプはより急（傾き比
  \(1.18\)--\(1.31\)）で閾値へより速く到達する。onset
  過大評価と非正規超過は同一機構の表裏。
\end{itemize}

\textbf{深い読み（接線コサイクルの有限時間 Lyapunov 率）。} 床近傍では
\(\delta_{t+1}=J_0\delta_t\) で率は
\(\ln\rho(J_0)=\ln|\mu_{\max}|\)（床近傍の初期成長率）。床を離れると実軌道上で
\(\delta_{t+1}=J(z_t)\delta_t\)、すなわち
\(\delta_t=\Phi(t,0)\delta_0\)、\(\Phi(t,0)=J(z_{t-1})\cdots J(z_0)\)（接線写像の時間順序積＝\textbf{接線コサイクル}）。したがって\textbf{実インフレーション率は、実軌道上の接線コサイクル
\(\Phi(t,0)=\prod_t J(z_t)\) の有限時間 Lyapunov 率}

\[\gamma_t=\frac1t\ln\frac{\|P_\perp\Phi(t,0)\delta_0\|}{\|P_\perp\delta_0\|},\qquad \log_{10}(H_\perp/H)\text{ の傾き}=\frac{2\gamma_t}{\ln10}\]

である（\(P_\perp\)＝床平面直交射影）。\(J_t\)
は非正規（\(0.53\)--\(0.76\)）で互いに非可換ゆえ
\(\Phi(t,0)\neq J_0^t\)、最大増幅方向が時間発展で回転しながら次の
\(J_t\) にさらに増幅される（unstable eigenmode → rotating amplified
direction → time-ordered exponential growth）。追試Aはこの \(\gamma_t\)
を実測ランプと一致（時変/実測
0.999／1.004／1.019）で確認した。\(\mu_{\max}\)（床の局所不安定性）と
\(\Phi(t,0)\)（実インフレの接線力学）は\textbf{二段階の力学}であり、\(\mu_{\max}\)
はその下界（\(t\to0^+\)・凍結極限）である。

\subsection{5.
二段構造の結論}\label{ux4e8cux6bb5ux69cbux9020ux306eux7d50ux8ad6}

\begin{itemize}
\tightlist
\item
  \textbf{導ける（形・機構・\(N\)
  依存・onset）}：インフレ曲線は床の線形不安定モードの指数増幅。主要予測（傾き
  \(2\lambda/\ln10\)・onset・\(N\) 依存）は床ヤコビアンの最大固有値
  \(\lambda=\ln|\mu_{\max}|\) が与える。
\item
  \textbf{論文5との統一（同一 \(\mu_{\max}\) の二つの顔）}：この
  \(\lambda=\ln|\mu_{\max}|\)
  は論文5のメキシカンハット有効ポテンシャルの頂上曲率そのものである（\(V''(0)=-\lambda\)、\(\lambda=\ln|\mu_{\max}|\)）。すなわち
  \textbf{SSB
  の不安定性（谷への転落を駆動する頂上の負曲率）とインフレーションの指数成長率は、同一の床ヤコビアン固有値
  \(\mu_{\max}\) で結ばれる}。論文5の静的 SSB
  像（不安定な対称極大からの自発選択）と本論文の動的増幅（指数急拡大）は、一つの
  \(\mu_{\max}\) の静的・動的な二側面である。\(N\) 増で
  \(\mu_{\max}\to1\)
  に近づくほど頂上は平坦化し（曲率減）、転落＝インフレも遅くなる（onset
  伸長）。増幅率が精度・シード深さに依存しない（＝力学に内在する）ことは、シードの底を100桁精度で検証した実験で確立済み------IC100（closure
  \(\sim10^{-102}\)・床 \(\sim10^{-200}\)）でも \(\gamma_\tau\)
  が64bit走行と8桁（\(N=8\)）一致し、床から約179桁を単一指数則で発火する（\texttt{precision\_seed\_dynamics\_isolation\_100digit\_20260902}、\(N=7,8\)）。
\item
  \textbf{精密化（非正規補正）}：正確な傾きは時変非正規 \(J\)
  の時間順序積の増幅で、床固有値に \(\sim18\)--\(31\%\)
  上乗せされる（\textbf{追試Aで直接確認}：時変積の傾きが実測ランプに一致し床凍結を18--31\%上回る）。閉形式のみ未導出。
\end{itemize}

\textbf{統一連鎖}（点火から着地まで一本）：

\[\rho(J_0)>1\ \Rightarrow\ \text{点火資格（論文2）}\ \to\ \ln\rho(J_0)\ \Rightarrow\ \text{床近傍の初期成長率}\ \to\ \gamma_{\rm FT}\big[\textstyle\prod_t J(z_t)\big]\ \Rightarrow\ \text{実ランプ成長率（追試A）}\ \to\ \frac{\ln(1/\text{seed})}{\gamma}\ \Rightarrow\ \text{onset}\ \to\ \text{非線形飽和}\ \Rightarrow\ \text{縮退真空（論文5）}.\]

同じ不安定床を「どう床から離れるか」で見ればインフレーション、「どこへ着地するか」で見れば
SSB------両者は別機構ではなく、一つの不安定床の二つの見方である。

これはメキシカンハット頂上（90°骨格の不安定な相対平衡）からの転がり落ちの定量記述であり、着地（縮退真空・\(\Delta=-2\pi/N\)）は論文5、床平面から傾く幾何は論文7が扱う。

\subsection{6.
主張分類・未決}\label{ux4e3bux5f35ux5206ux985eux672aux6c7a}

\begin{itemize}
\tightlist
\item
  分類：導出済み帰結（線形不安定式）＋非正規補正の実測（精密化は仮説）。判定：保持。
\item
  未決：非正規時変積の増幅率の\textbf{閉形式}は未導出（機構は追試Aで確定、増幅率の解析式のみ今後の課題）。\(N=40\)
  の床ヤコビアン（\(2M=1560\) 次元）は重く未計算------小
  \(N\)（6,8,10）で二段構造は確立。
\end{itemize}

\subsection{関連研究（構造的一致・導出元でない）}\label{ux95a2ux9023ux7814ux7a76ux69cbux9020ux7684ux4e00ux81f4ux5c0eux51faux5143ux3067ux306aux3044}

\begin{itemize}
\tightlist
\item
  非正規行列の一時的増幅・擬スペクトル（Trefethen \& Embree
  2005）：ランプが床固有値を超える分は非正規 transient
  growth／時間順序積の増幅。
\item
  自己無撞着インフレv2［1］：本論文はその急拡大機構の解析導出（精密化）。
\end{itemize}

\subsection{参考文献}\label{ux53c2ux8003ux6587ux732e}

\begin{enumerate}
\def\labelenumi{\arabic{enumi}.}
\tightlist
\item
  木原範昭「自己無撞着な関係波閉鎖系におけるインフレーション的急拡大の機構」Version
  DOI 10.5281/zenodo.22176949 / Concept DOI 10.5281/zenodo.22112008.
\item
  L. N. Trefethen and M. Embree, \emph{Spectra and Pseudospectra},
  Princeton (2005).
\end{enumerate}

\subsection{再現}\label{ux518dux73fe}

検証スクリプト・図・報告は
\texttt{位相ロック全N確認\_相対平衡\_20260912/インフレ曲線\_線形不安定解析\_20260912/}（\texttt{analyze\_inflation\_linear\_instability\_20260912.py}＝検証A〜E［床\(J\)固有値／非正規度／固定床\(\delta\)発展／実測ランプ／onset］、\texttt{make\_figure\_inflation.py}、\texttt{fig\_inflation\_linear\_instability.png}、\texttt{REPORT.md}、\texttt{解析ログ\_20260912.log}、\texttt{SHA256SUMS.txt}、\texttt{run\_all.sh}）。\textbf{非正規補正の機構の直接確認（追試A）}は
\texttt{位相ロック全N確認\_相対平衡\_20260912/ランプ直接Lyapunov\_非正規時変積\_20260912/}（\texttt{analyze\_ramp\_lyapunov\_20260912.py}＝床凍結／時変積／実測の
\(H_\perp\)
傾き比較、\texttt{results/ramp\_lyapunov\_summary.csv}、\texttt{README.md}／\texttt{REPORT.md}／\texttt{SHA256SUMS.txt}／\texttt{run\_all.sh}）。入力は既存
den=N npz（500 step、make\_parent
床）。走行・物理変更なし。力学写像は正本
\texttt{run\_N3\_N40\_stage123\_v1.py}（SHA256
\texttt{1abf2353fee2e4f56f05e7a6f149fd086885136beb61ab571b48a56b09691567}）と同一。

\end{document}
